\documentclass[11pt]{article}
\usepackage{fontspec}
\setmainfont{DejaVu Serif}
\setmonofont{DejaVu Sans Mono}[Scale=MatchLowercase]
\usepackage[a4paper,margin=25mm]{geometry}
\usepackage{amsmath,amssymb,booktabs,tabularx,hyperref,microtype}
\newcolumntype{Y}{>{\raggedright\arraybackslash}X}
\setcounter{secnumdepth}{0}
\hypersetup{colorlinks=true,linkcolor=blue,urlcolor=blue}
\title{A theorem-licensed typestate machine for exact rotation algebra}
\author{John Curley\\Independent Researcher, SPU-13 Project\\Wellington, New Zealand}
\date{10 August 2026 (working draft toward version 1.2; published version is 1.1)}
\begin{document}
\maketitle
\noindent\textbf{AI disclosure.} AI tools assisted with drafting, repository evidence synthesis, and editorial revision. The human author reviewed the manuscript, selected its claims and scope, and is responsible for the submitted record.

\section{Abstract}

[THEOREM] This paper presents a typestate construction for exact rotation algebra whose transitions are licensed by a doubling theorem. The construction separates fresh values from values known to remain in one catalog domain, and rejects transitions when the theorem's validity domain is not established. [RTL] The construction is instantiated in the ROTC and IROTC paths and is connected to a bounded tensegrity admission guard. [SILICON] Selected ROTC, IROTC, and tensegrity paths have board evidence, while the complete catalog surfaces remain a mixture of RTL and board-scoped results. The contribution is the correspondence between an algebraic domain and an enforced transition relation; it is not a claim about project-wide evidence or arbitrary arithmetic values.

\section{SPU-13 context}

[RTL] SPU-13 is the engineering host for these case studies: ROTC supplies integer-coordinate rotation paths, IROTC supplies exact Z[φ] catalog rotations, and the tensegrity guard supplies a bounded admission state machine. They share an exact-arithmetic hardware setting, but the paper does not claim that the project as a whole is covered by the typestate result. The examples are chosen to show the same transition discipline at arithmetic, catalog, and system boundaries.

\section{1. Scope and contribution}

[THEOREM] The central result is a four-state typestate lattice for registers holding Z[φ] pairs: \path{UNTAGGED}, \path{FRESH}, \path{MAIN}, and \path{CONJ}. \path{MAIN} and \path{CONJ} mean that the register is known to lie in one of two catalog domains; they do not mean that an arbitrary value has been proved to be a desired application result.

[THEOREM] Within one A₅ catalog, the doubling theorem licenses the transition from \path{FRESH} to the catalog's typed state and licenses subsequent operations that preserve that domain. The conjugate catalog carries the Galois-dual domain. A mixed main/conjugate product is outside the theorem's domain and is therefore rejected rather than silently truncated.

[THEOREM] \path{PCHIRAL} carries the ring automorphism between the two catalog domains. \path{SCALE2} reconditions a register to \path{FRESH}. Operations outside the A₅ theorem domain, including the integer-matrix ROTC classes, either preserve only the state they are specified to preserve or demote the tag to \path{UNTAGGED}.

[RTL] The state machine is represented in the IROTC integration as a two-bit tag per QR register, with dispatch faults for untagged use, bad catalog index, and catalog mixing. Faulting operations hold the destination and its tag.

[THEOREM] The result is a transition-soundness statement: every accepted transition is licensed by the stated precondition and theorem domain. It does not imply that the datapath computes the intended value, that a testbench exhausts all values, or that a document claim is supported.

\section{2. Formal model}

\subsection{2.1 State space}

[THEOREM] Let the register typestate be

\(\tau \in \{U,F,M,C\}\)

for \path{UNTAGGED}, \path{FRESH}, \path{MAIN}, and \path{CONJ}. Let \path{D_M} and \path{D_C} denote the two catalog domains and let \path{φ} denote the ring automorphism used by the conjugate construction. The state invariant is a judgment of the form \path{Γ ⊢ (r, τ)}, where \path{M} and \path{C} carry a domain witness and \path{U} carries none.

[THEOREM] The principal transition rules are:

\begin{center}

\scriptsize

\sloppy

\begin{tabularx}{\linewidth}{@{}YYY@{}}

\toprule

Operation & Precondition & Result \\

\midrule

raw load & none & \path{U} \\

\path{SCALE2} & \path{U}, \path{M}, or \path{C} & \path{F} \\

main-catalog IROTC & \path{F} or \path{M} & \path{M} \\

conjugate-catalog IROTC & \path{F} or \path{C} & \path{C} \\

\path{PCHIRAL} & \path{M} or \path{C} & the opposite typed domain \\

mixed catalog operation & \path{M} and \path{C} & reject with \path{CATMIX} \\

octahedral/integer ROTC & operation-specific & preserve or demote as specified \\

\bottomrule

\end{tabularx}

\end{center}

[THEOREM] The induction argument is by transition cases. Raw writes establish no catalog fact; reconditioning establishes the \path{FRESH} precondition; a same-catalog step consumes and returns the same domain witness; \path{PCHIRAL} applies the automorphism; and a mixed step has no licensed conclusion. Thus a rejected transition is evidence that the machine has reached the boundary of the theorem's stated domain, not evidence that the underlying value is invalid.

\subsection{2.2 What the theorem does not establish}

[THEOREM] The judgment is deliberately narrower than a value proof. It does not establish bounds, expected test-vector outputs, transport framing, clock configuration, placement quality, or the absence of faults outside the transition relation.

\section{3. Case studies}

\subsection{3.1 ROTC}

[RTL] The tagged ROTC path models \path{CLEAN}, \path{PENDING}, and explicit \path{MISALIGNED}, \path{OVERFLOW}, and \path{INEXACT} fault states. Its acceptance harness provides 9/9 acceptance tests in \path{spu13_rotor_core_tagged_tb.v}, and the tagged core is distinct from the original TDM baseline.

[SILICON] ROTC angles 0–5 have board evidence in \href{https://github.com/pinnacletechnologysolutionsltd/SPU/blob/v1.1-typestate/docs/hardware_evidence.md#32g-rotc-0-5-silicon-probe}{hardware\_evidence.md §3.2g (ROTC 0–5)}.

[RTL] ROTC angles 6–35 are covered by RTL/trace evidence only; they are not reported here as board results. The integer permutation classes are therefore an example of a transition surface whose algebraic classification and board status must remain separate.

\subsection{3.2 IROTC}

[THEOREM] IROTC is the direct case study for the lattice: main-catalog chains remain in \path{MAIN}, conjugate chains remain in \path{CONJ}, \path{PCHIRAL} changes the catalog witness, and mixed chains have no theorem-licensed result.

[RTL] The RTL engine covers the generated 60-entry main catalog and its conjugate catalog, including the typestate dispatch and poison-hold behavior. The full 60×2 catalog is reported as RTL-only in this paper.

[SILICON] The IROTC probe vectors and the \path{BADIDX}/\path{UNTAGGED}/\path{CATMIX} fault matrix are supported by \href{https://github.com/pinnacletechnologysolutionsltd/SPU/blob/v1.1-typestate/docs/hardware_evidence.md#32k-irotc-icosahedral-rotation-engine-silicon-probe}{hardware\_evidence.md §3.2k}.

[SILICON] SPI core integration and the conjugate-catalog cases are supported by \href{https://github.com/pinnacletechnologysolutionsltd/SPU/blob/v1.1-typestate/docs/hardware_evidence.md#32k1-irotc-spi-core-integration--conjugate-catalog-silicon-proof}{hardware\_evidence.md §3.2k.1}.

\subsection{3.3 Tensegrity admission}

[RTL] The tensegrity guard demonstrates the same separation at a system boundary: a bounded admission decision carries explicit state, fault, and recovery transitions, while its geometry and equilibrium predicates remain separate computations.

[SILICON] The seven frozen admission fixtures, including the equilibrium fixture, are supported by \href{https://github.com/pinnacletechnologysolutionsltd/SPU/blob/v1.1-typestate/docs/hardware_evidence.md#32l-wukong-tensegrity-admission-guard-silicon-probe}{hardware\_evidence.md §3.2l — V:7 final admission tranche}.

[SILICON] The four-act transport and recovery sequence is supported by \href{https://github.com/pinnacletechnologysolutionsltd/SPU/blob/v1.1-typestate/docs/hardware_evidence.md#32l1-tensegritylink-four-act-proof-on-the-karatsuba-candidate}{hardware\_evidence.md §3.2l.1}.

\section{4. Harness and coverage boundary}

[RTL] The repository's state-machine document lists typestate or explicit state descriptions for ROTC, SOM/BMU, BTU, Padé evaluation, Lucas MAC, and batch inversion. Its completed typestate harness is the ROTC item; the other entries describe planned or partial harness work and must not be counted as completed formal coverage.

\begin{center}

\scriptsize

\sloppy

\begin{tabularx}{\linewidth}{@{}YYY@{}}

\toprule

Subsystem & State-machine material & Coverage represented here \\

\midrule

ROTC & Completed tagged-state harness & RTL harness; angles 0–5 additionally have SILICON evidence \\

SOM/BMU & State description and planned harness & Testbench/oracle material only; no completed typestate harness claimed \\

BTU & State description and planned harness & Testbench/oracle material only; no completed typestate harness claimed \\

Padé evaluator & State description and planned harness & Testbench/oracle material only; no completed typestate harness claimed \\

Lucas MAC & State description and partial checks & Testbench/oracle material only for this paper \\

Batch inversion & State description and invariants & Testbench/oracle material only; no completed typestate harness claimed \\

IROTC & Integrated typestate tags and fault transitions & RTL engine/integration; selected cases SILICON, full catalog RTL-only \\

Tensegrity guard & Explicit admission/fault/recovery states & RTL guard; seven fixtures and transport also SILICON \\

\bottomrule

\end{tabularx}

\end{center}

[RTL] The distinction in this table is intentional: a bench that checks outputs or protocol frames is not counted as a typestate harness, and a state-machine description is not counted as an executed proof artifact.

\section{5. Negative results}

[RTL] The planned SOM/BMU, BTU, Padé, Lucas MAC, and batch-inversion harnesses are specified in the state-machine document but are not presented as built typestate harnesses here. This is a negative result about the scope of the current artifact.

[THEOREM] Several fault classes cannot be expressed by this machine: a wrong arithmetic value produced while all tags are valid, an omitted clock divider, a placement-dependent timing failure, a malformed evidence statement, and a transport or physical failure outside the modeled transition relation. Those failures require independent value checks, build evidence, or bench evidence.

[THEOREM] The theorem domain and RTL enforcement are also different claims. The theorem licenses one A₅ domain and its conjugate transformation; RTL can store and check tags, but that alone does not prove every theorem assumption or every datapath result. The full generated 60×2 catalog therefore remains an RTL claim in this paper, not a blanket board claim.

\section{6. Limitations}

[THEOREM] Typestate constrains transitions, not values. A well-typed execution may still compute the wrong value, and the machine does not police claims, documentation, papers, or status reports.

[OBSERVED] The audit of this paper is deliberately separate from the draft. The five defects found on 2026-08-09/10 were: the aliased loop variable that ran one vector instead of six; the coverage report that read as complete; the critical- path misattribution; the RPLU2PADE clock omission; and the seven unbacked bullets. The typestate machine would have caught none of them.

[THEOREM] The formal result is therefore a bounded transition result. It is not a substitute for independent vector coverage, synthesis and timing inspection, configuration accounting, or a claim-by-claim evidence audit.

\section{7. Evidence ledger for this draft}

Every technical claim in the paper is assigned one tier. The tier is about the kind of support available for that claim, not a ranking of the theorem.

\begin{center}

\scriptsize

\sloppy

\begin{tabularx}{\linewidth}{@{}YYYYY@{}}

\toprule

ID & Claim summary & Tier & Source & Body mapping \\

\midrule

T1 & Four-state lattice and domain witnesses & THEOREM & §2.1 & §1 lines 37–41; §2.1 lines 68–77 \\

T2 & Same-catalog doubling transition & THEOREM & §2.1 & §1 lines 43–47; §2.1 lines 79–97 \\

T3 & \path{PCHIRAL}, \path{SCALE2}, and integer-ROTC boundary rules & THEOREM & §2.1 & §1 lines 49–53; §2.1 lines 79–97 \\

T4 & Transition soundness is narrower than value proof & THEOREM & §1, §2.2 & §1 lines 59–62; §2.2 lines 101–104 \\

T5 & IROTC main/conjugate transition interpretation & THEOREM & §3.2 & §3.2 lines 125–127 \\

T6 & Unexpressible fault classes and theorem/RTL gap & THEOREM & §5 & §5 lines 182–192 \\

R1 & Tagged ROTC states and tagged-state harness & RTL & \path{spu13_rotor_core_tagged_tb.v} (9/9 acceptance tests, plus 2000 randomized REDUCE cases) & §3.1 lines 110–113 \\

R2 & ROTC angles 6–35 are RTL/trace-only here & RTL & \path{test_rotc_vm_rtl_trace.py} (336 checks over 42 generated cases) & §3.1 lines 118–121 \\

R3 & IROTC generated 60×2 surface and fault transitions & RTL & \path{spu13_irotc_engine_tb.v} (120 oracle golden cases; fixed 12-clock latency) & §3.2 lines 129–131 \\

R4 & Tensegrity bounded state/fault/recovery model & RTL & \path{spu13_tensegrity_guard_tb.v} (9/9 cases); tensegrity suite total 50 PASS (suite total, not a check count) & §3.3 lines 141–144 \\

R5 & Harness-versus-bench coverage boundary & RTL & \path{STATE_MACHINE_HARNESS.md} §3 (8 subsystem entries; 1 completed harness) & §4 lines 154–173 \\

R6 & SPU-13 case-study context & RTL & \path{spu13_core_rotc_opcode_tb.v} (41 lane checks + 2 poison faults) and \path{spu13_core_irotc_opcode_tb.v} (25 checks) & §SPU-13 context lines 27–33 \\

S1 & ROTC angles 0–5 & SILICON & \path{hardware_evidence.md} §3.2g (ROTC 0–5) & §3.1 lines 115–116 \\

S2 & IROTC probe and fault matrix & SILICON & \path{hardware_evidence.md} §3.2k & §3.2 lines 133–134 \\

S3 & IROTC SPI integration and conjugate catalog & SILICON & \path{hardware_evidence.md} §3.2k.1 & §3.2 lines 136–137 \\

S4 & Tensegrity seven admission fixtures & SILICON & \path{hardware_evidence.md} §3.2l & §3.3 lines 146–147 \\

S5 & Tensegrity four-act transport and recovery & SILICON & \path{hardware_evidence.md} §3.2l.1 & §3.3 lines 149–150 \\

O1 & Five repository defects observed 2026-08-09/10; typestate caught none & OBSERVED & §6 & §6 lines 200–204 \\

\bottomrule

\end{tabularx}

\end{center}

The repository sections \path{§3.1–§3.2c} are only partial backing for the broader project context: they do not supply the date and SHA-256 shape required for a silicon claim in this paper, so none of them is used as a SILICON source here.

\section{References}

All repository citations below are pinned to the \path{v1.1-typestate} release tag.

\begin{enumerate}

\item \href{https://github.com/pinnacletechnologysolutionsltd/SPU/blob/v1.1-typestate/docs/STATE_MACHINE_HARNESS.md#3-subsystem-state-machines}{STATE\_MACHINE\_HARNESS.md §3}.

\item \href{https://github.com/pinnacletechnologysolutionsltd/SPU/blob/v1.1-typestate/docs/IROTC_SPEC.md}{IROTC\_SPEC.md}.

\item \href{https://github.com/pinnacletechnologysolutionsltd/SPU/blob/v1.1-typestate/docs/hardware_evidence.md#32g-rotc-0-5-silicon-probe}{hardware\_evidence.md §3.2g (ROTC 0–5)}, \href{https://github.com/pinnacletechnologysolutionsltd/SPU/blob/v1.1-typestate/docs/hardware_evidence.md#32k-irotc-icosahedral-rotation-engine-silicon-probe}{§§3.2k, 3.2k.1, 3.2l, 3.2l.1}.

\item \href{https://github.com/pinnacletechnologysolutionsltd/SPU/blob/v1.1-typestate/docs/ROTC_EXPONENT_STATE_MACHINE.md}{ROTC\_EXPONENT\_STATE\_MACHINE.md}.

\end{enumerate}
\end{document}
